\documentclass{article}
\usepackage[T1]{fontenc}
\usepackage{lmodern}
\usepackage{graphicx}
\usepackage{amsmath,amssymb}
\usepackage[a4paper,margin=2cm]{geometry}
\usepackage[absolute,overlay]{textpos}
\setlength{\TPHorizModule}{1mm}
\setlength{\TPVertModule}{1mm}
\usepackage[indonesian]{babel}
\usepackage{hyperref}
\setlength{\emergencystretch}{2em}
% unit: Halaman 5

\begin{document}

% === Halaman 5 (python/native) ===

Penurunan Rumus Enkripsi dan Dekripsi RSA

• Teori yang digunakan: Teorema Euler    a(n)  1 (mod n) 

• Syarat: 

1. a harus relatif prima terhadap n , yaitu PBB(a, n) = 1          (PBB = gcd)

2. (n) = Toitent Euler = fungsi yang menentukan berapa  banyak  dari bilangan-

bilangan 1, 2, 3, …, n yang relatif prima terhadap n.


Contoh: (20) = 8, sebab terdapat 8 buah  yang relatif  prima dengan 20, yaitu  

1, 3, 7, 9, 11, 13, 17, 19. 


      Jika n = pq adalah bilangan komposit dengan p dan q  prima, maka

             (n) = (p) (q) = (p – 1)(q – 1).

5

\end{document}
